\documentclass[11pt]{article}

\usepackage[utf8]{inputenc}
\usepackage[T1]{fontenc}
\usepackage{amsmath}
\usepackage{graphicx}
\usepackage{booktabs}
\usepackage{natbib}
\usepackage{hyperref}
\usepackage[margin=1in]{geometry}
\usepackage{xcolor}
\usepackage{lineno}

\hypersetup{
    colorlinks=true,
    linkcolor=blue,
    citecolor=blue,
    urlcolor=blue
}

\title{Disseminating Cells in Human Oral Tumours Acquire an EMT Cancer Stem Cell State That Is Predictive of Metastasis}

\author{
Author One$^{1,*}$, Author Two$^{1}$, Author Three$^{2}$, Author Four$^{1}$, Author Five$^{3}$ \\[6pt]
\small $^{1}$Centre for Tumour Biology, Barts Cancer Institute, Queen Mary University of London, UK \\
\small $^{2}$Department of Pathology, University of Example, City, Country \\
\small $^{3}$Department of Oral Surgery, University of Example, City, Country \\
\small $^{*}$Corresponding author: \texttt{author.one@example.ac.uk}
}

\date{}

\begin{document}

\maketitle

\begin{abstract}
The epithelial-to-mesenchymal transition (EMT) is widely proposed to drive metastatic dissemination, yet direct observation of EMT-derived cancer cells leaving human tumours has not been achieved. A fundamental challenge is that cancer cells undergoing EMT downregulate epithelial markers, becoming indistinguishable from surrounding mesenchymal stromal cells. Here we exploit the selective retention of EpCAM by a subset of EMT cancer stem cells (CSCs) to identify disseminating cells in the stroma of human oral squamous cell carcinomas. Using automated four-colour immunofluorescence imaging across entire histopathological slides from 84 tumour specimens (over 12,000 imaging fields), we show that EpCAM$^{+}$Vimentin$^{+}$CD24$^{+}$ triple-positive cells are enriched to 41\% in an EMT-stem sub-line versus 2.1\% in the parental CA1 line, and are absent in normal keratinocytes and cancer-associated fibroblasts. Single EpCAM$^{+}$Vim$^{+}$CD24$^{+}$ cells disseminating beyond the tumour body are significantly enriched in the stroma of metastatic compared to non-metastatic specimens, whereas EpCAM$^{+}$Vim$^{+}$CD24$^{-}$ cells show no such enrichment. Replacing EpCAM with pan-keratin eliminates the signal, confirming that keratins are lost during EMT while EpCAM is retained. An artificial neural network trained on the three-marker staining profile predicts metastasis with 87--89\% cross-validated accuracy, maintained on an independent validation batch. These findings provide the first direct observation of EMT CSCs disseminating from human tumours \textit{in vivo} and establish a predictive link between disseminating EMT CSCs and metastatic outcome.
\end{abstract}

\noindent\textbf{Keywords:} epithelial-mesenchymal transition, cancer stem cells, metastasis, oral squamous cell carcinoma, EpCAM, dissemination

\section{Introduction}

Metastasis is the leading cause of cancer-related mortality, yet the cellular mechanisms by which tumour cells detach from the primary mass and disseminate to distant sites remain incompletely understood in human cancers \citep{Hanahan2011,Lambert2017}. The epithelial-to-mesenchymal transition (EMT) has emerged as a central paradigm for understanding metastatic dissemination \citep{Thiery2002,Nieto2016}. During EMT, epithelial cancer cells acquire mesenchymal characteristics---including enhanced motility, invasiveness, and resistance to apoptosis---enabling them to leave the primary tumour, invade surrounding tissues, and enter the vasculature \citep{Yang2004,Mani2008}.

However, the EMT model of metastasis has been built almost entirely from mouse models and cell line studies \citep{Rhim2012,Fischer2015,Zheng2015}. Directly observing EMT-driven dissemination in human tumour specimens presents a fundamental technical challenge: as cancer cells undergo EMT, they downregulate epithelial markers such as keratins and E-cadherin, becoming morphologically and immunophenotypically indistinguishable from the mesenchymal stromal cells that surround the tumour \citep{Kalluri2009,Ye2015}. Consequently, once an EMT cancer cell has left the tumour body and entered the stroma, it becomes effectively invisible to standard histological and immunohistochemical analysis.

Previous attempts to characterize EMT in human tumours have relied on detecting cells at the tumour-stroma boundary that co-express epithelial and mesenchymal markers, capturing only the earliest stages of EMT in cells still physically attached to the tumour mass \citep{Brabletz2001,Savagner2010}. Single-cell RNA sequencing has identified partial EMT gene expression signatures correlated with metastatic potential \citep{Puram2017,Pastushenko2018}, but this approach lacks spatial resolution and cannot determine whether cells have physically disseminated from the tumour.

A potential solution comes from the observation that the epithelial cell adhesion molecule EpCAM is selectively retained by a subset of cancer cells that have undergone EMT and acquired stem cell properties \citep{Grosse-Wilde2015,Bierie2017}. Unlike pan-keratins, which are lost during EMT, EpCAM remains on the surface of these EMT CSCs, providing a unique tumour lineage tracer in the stromal compartment where mesenchymal cells predominate. Furthermore, CD24 has been identified as a marker of EMT CSCs with enhanced plasticity---the ability to undergo the reverse mesenchymal-to-epithelial transition (MET) required for metastatic colonization at distant sites \citep{Al-Hajj2003,Grosse-Wilde2015}.

Here we combine EpCAM, Vimentin (a mesenchymal marker upregulated during EMT), and CD24 in a four-colour immunofluorescent staining protocol applied to entire histopathological slides from 84 human oral squamous cell carcinoma (OSCC) specimens. Using automated imaging and computational analysis, we identify single disseminating EMT CSCs in the tumour stroma and demonstrate their predictive association with metastatic disease.

\section{Results}

\subsection{EpCAM$^{+}$Vim$^{+}$CD24$^{+}$ identifies EMT cancer stem cells \textit{in vitro}}

We first validated the three-marker staining panel using the CA1 oral squamous cell carcinoma cell line and its derivative EMT-stem sub-line, which has been previously characterized as possessing EMT and cancer stem cell properties \citep{Grosse-Wilde2015}. Four-colour immunofluorescent staining was performed using EpCAM (yellow), Vimentin (red), CD24 (green), and DAPI (blue) as a nuclear counterstain.

\begin{table}[ht]
\centering
\caption{Triple-positive cell frequencies in cell line validation experiments.}
\label{tab:celllines}
\begin{tabular}{lcc}
\toprule
Cell Population & EpCAM$^{+}$Vim$^{+}$CD24$^{+}$ (\%) & Pan-keratin$^{+}$Vim$^{+}$CD24$^{+}$ (\%) \\
\midrule
CA1 parental line & 2.1 & Very low \\
EMT-stem sub-line & 41.0 & Very low (no enrichment) \\
Normal keratinocytes & 0 (absent) & 0 (absent) \\
Cancer-associated fibroblasts & 0 (absent) & 0 (absent) \\
\bottomrule
\end{tabular}
\end{table}

The EpCAM$^{+}$Vim$^{+}$CD24$^{+}$ triple-positive population was enriched from 2.1\% in the parental CA1 line to 41.0\% in the EMT-stem sub-line ($p < 0.05$, two-tailed Student's $t$-test; Table~\ref{tab:celllines}). Critically, triple-positive cells were completely absent in normal keratinocytes and cancer-associated fibroblasts, confirming that the staining profile is specific to cancer cells that have undergone EMT while retaining stem cell plasticity.

A parallel staining panel substituting pan-keratin for EpCAM showed no enrichment of triple-positive cells in the EMT-stem sub-line (Table~\ref{tab:celllines}), confirming that pan-keratins are lost during EMT whereas EpCAM is retained. This result validates EpCAM as the appropriate tumour lineage marker for detecting disseminating EMT cells in the stromal compartment.

\subsection{Automated whole-slide imaging and tumour-stroma delineation}

To detect disseminating cells in human tumour specimens, we developed an automated imaging and analysis pipeline applied to entire histopathological slides (Table~\ref{tab:imaging}).

\begin{table}[ht]
\centering
\caption{Image analysis parameters for tumour specimen analysis.}
\label{tab:imaging}
\begin{tabular}{ll}
\toprule
Parameter & Value \\
\midrule
Imaging mode & Automated 4-colour immunofluorescence \\
Coverage & Entire histopathological slides \\
Total imaging fields & $>$12,000 \\
Total specimens & 84 \\
Segmentation output & Per-cell marker intensities \\
Tumour-stroma boundary & EpCAM density map \\
Cell identification & Individual cell level \\
\bottomrule
\end{tabular}
\end{table}

An EpCAM density map was computed across each slide to objectively delineate the tumour body (high EpCAM density) from the surrounding stroma (low EpCAM density). Disseminating cells were defined as single EpCAM$^{+}$ cells located in the stromal compartment, having physically separated from the main tumour mass.

\subsection{EpCAM$^{+}$Vim$^{+}$CD24$^{+}$ disseminating cells are enriched in metastatic tumour stroma}

Eighty-four human OSCC specimens, stratified by metastatic status, were analyzed in two independent batches. Four-colour staining (EpCAM, Vimentin, CD24, DAPI) was applied, and the frequency of triple-positive cells in the stromal compartment was quantified.

\begin{table}[ht]
\centering
\caption{Enrichment of marker combinations in metastatic versus non-metastatic tumour stroma.}
\label{tab:enrichment}
\begin{tabular}{lll}
\toprule
Marker Combination & Enrichment in Metastatic Stroma & Clinical Relevance \\
\midrule
EpCAM$^{+}$Vim$^{+}$CD24$^{+}$ & Significant ($p < 0.05$) & Correlates with metastasis \\
EpCAM$^{+}$Vim$^{+}$CD24$^{-}$ & No significant enrichment & Not predictive \\
Pan-keratin$^{+}$Vim$^{+}$CD24$^{+}$ & No enrichment & Not useful \\
\bottomrule
\end{tabular}
\end{table}

Single EpCAM$^{+}$Vim$^{+}$CD24$^{+}$ cells disseminating beyond the tumour body were significantly enriched in the stroma of metastatic specimens compared to non-metastatic specimens ($p < 0.05$; Table~\ref{tab:enrichment}). Critically, EpCAM$^{+}$Vim$^{+}$CD24$^{-}$ cells in the stroma did \textit{not} correlate with metastasis, demonstrating that the CD24$^{+}$ plasticity marker is specifically required for identifying disseminating cells predictive of metastatic outcome. Furthermore, staining with the pan-keratin control panel showed no enrichment, confirming that keratins are lost in these disseminating cells.

\begin{table}[ht]
\centering
\caption{Staining quantification across specimen types in the stromal compartment.}
\label{tab:staining}
\begin{tabular}{lcccc}
\toprule
Specimen Type & EpCAM$^{+}$ & Vim$^{+}$ & CD24$^{+}$ & Triple-positive \\
\midrule
Normal epithelium & Minimal & High & Variable & Absent/rare \\
Non-metastatic OSCC & Low & High & Variable & Low frequency \\
Metastatic OSCC & Elevated & High & Variable & Significantly enriched \\
\bottomrule
\end{tabular}
\end{table}

\subsection{Machine learning predicts metastasis from three-marker profile}

To quantify the discriminative power of the three-marker staining profile, we trained an artificial neural network (ANN) on grey-level intensities of EpCAM, Vimentin, and CD24 in imaging tiles from tumour batch 1 (Table~\ref{tab:ml}).

\begin{table}[ht]
\centering
\caption{Machine learning model specifications and performance.}
\label{tab:ml}
\begin{tabular}{ll}
\toprule
Component & Specification \\
\midrule
Architecture & Artificial neural network \\
Input features & Grey-level intensities (3 channels) per tile \\
Labels & Metastatic vs.\ non-metastatic tumour origin \\
Cross-validation & 10-fold (batch 1) \\
Cross-validated accuracy & 87--89\% \\
External validation & Independent batch 2 \\
External validation accuracy & High (maintained) \\
\bottomrule
\end{tabular}
\end{table}

The ANN achieved 87--89\% cross-validated accuracy on batch 1 using the EpCAM-based panel. When the same approach was applied using the pan-keratin control panel (pan-keratin, Vimentin, CD24), classification accuracy was substantially lower, further confirming the specificity of EpCAM retention in the EMT CSC state.

\begin{table}[ht]
\centering
\caption{Independent validation of the machine learning classifier.}
\label{tab:validation}
\begin{tabular}{llll}
\toprule
Evaluation & Training Set & Test Set & Accuracy \\
\midrule
Within-batch (10-fold CV) & Batch 1 & Batch 1 (held-out folds) & 87--89\% \\
Cross-batch validation & Batch 1 & Batch 2 (independent) & High (maintained) \\
\bottomrule
\end{tabular}
\end{table}

The model trained on batch 1 was applied to batch 2 as an independent test set, where high classification accuracy was maintained (Table~\ref{tab:validation}), demonstrating generalizability across independently collected specimen cohorts.

\subsection{Marker biology underpinning the triple-positive state}

The biological basis for the discriminative power of the EpCAM$^{+}$Vim$^{+}$CD24$^{+}$ profile lies in the distinct expression dynamics of each marker during EMT (Table~\ref{tab:markers}).

\begin{table}[ht]
\centering
\caption{Marker expression patterns in relevant cell populations.}
\label{tab:markers}
\begin{tabular}{lllll}
\toprule
Marker & Role & EMT CSCs & Stroma & Lost in EMT? \\
\midrule
EpCAM & Epithelial adhesion & Retained & Absent & No (in plastic subset) \\
CD24 & CSC plasticity & Retained & Variable & No (in MET-competent) \\
Vimentin & Mesenchymal marker & Upregulated & High & N/A (mesenchymal) \\
Pan-keratin & Epithelial filaments & Lost & Absent & Yes \\
\bottomrule
\end{tabular}
\end{table}

EpCAM is an epithelial cell adhesion molecule that is specifically retained on the surface of cancer cells that undergo EMT while maintaining the plasticity to later undergo MET \citep{Grosse-Wilde2015}. This retention makes EpCAM uniquely suited as a tumour lineage tracer in the stromal compartment, where it is absent from resident fibroblasts and other mesenchymal cells. CD24 marks the subset of EMT cells with enhanced MET capability, which is believed to be essential for metastatic colonization \citep{Chaffer2013}. Vimentin, upregulated during EMT, serves as the mesenchymal marker confirming that the EpCAM$^{+}$CD24$^{+}$ cells have indeed undergone EMT rather than simply being detached epithelial cells.

\section{Discussion}

\subsection{First direct observation of disseminating EMT CSCs in human tumours}

Our findings provide the first direct observation of cancer cells that have undergone EMT and physically disseminated from the primary tumour mass in human specimens. Previous studies of EMT in human tumours were limited to cells at the tumour margin still attached to the tumour body \citep{Brabletz2001}, or to transcriptomic signatures lacking spatial resolution \citep{Puram2017}. Mouse lineage tracing studies have demonstrated EMT-driven dissemination \citep{Rhim2012}, but whether these findings translate to human cancer has remained uncertain. By exploiting the selective retention of EpCAM during EMT, we can identify tumour-derived cells that have entered the stromal compartment---a capability that is impossible with pan-keratin-based approaches, as keratins are lost during EMT.

\subsection{CD24 as a specificity marker for metastasis-associated dissemination}

A critical finding is that the CD24$^{+}$ plasticity marker is required for the metastatic association. EpCAM$^{+}$Vim$^{+}$ cells without CD24 expression are present in the stroma but do not correlate with metastatic outcome. This distinction is consistent with the emerging model that not all EMT states are equivalent: only cells in a hybrid EMT state retaining specific stem cell markers possess the plasticity to undergo MET at distant sites, a step considered essential for metastatic colonization \citep{Pastushenko2018,Jolly2015,Kroger2019}.

The requirement for CD24 also provides a mechanistic explanation for why bulk EMT markers alone have been poor predictors of metastasis in clinical studies \citep{Taube2010}. Our data suggest that the relevant prognostic information resides specifically in the subset of EMT cells retaining plasticity markers, not in the total population of cells exhibiting mesenchymal features.

\subsection{Implications for metastasis prediction and therapeutic targeting}

The 87--89\% cross-validated accuracy of the ANN classifier, maintained on independent validation, suggests clinical utility for metastasis prediction in oral cancer. Current prognostic tools for OSCC rely primarily on tumour staging and histological grading, which have limited accuracy for predicting nodal metastasis \citep{Woolgar2006,Almangush2015}. A quantitative immunofluorescence-based assay targeting the EpCAM$^{+}$Vim$^{+}$CD24$^{+}$ profile could provide additional prognostic information to guide treatment decisions, particularly regarding the extent of lymph node dissection.

Furthermore, the identification of a specific disseminating cell state suggests therapeutic opportunities. Targeting the EpCAM$^{+}$CD24$^{+}$ population through antibody-drug conjugates or bispecific antibodies could potentially intercept the metastatic cascade at the dissemination stage \citep{Soysal2015,Gires2020}. The retention of EpCAM on these cells, despite loss of other epithelial markers, provides a molecular handle for targeted intervention that would spare normal stromal cells.

\subsection{Comparison with prior approaches}

Our approach differs fundamentally from prior studies of EMT in human tumours (Table~\ref{tab:comparison}).

\begin{table}[ht]
\centering
\caption{Comparison of approaches for detecting EMT in human tumours.}
\label{tab:comparison}
\begin{tabular}{lll}
\toprule
Study Type & EMT Detection & Limitation \\
\midrule
Previous \textit{in vivo} & Cells at tumour margin & Cannot track detached cells \\
scRNAseq (partial EMT) & Gene expression & No spatial resolution \\
Mouse models & Lineage tracing & May not translate to human \\
This study & Single disseminating cells & First in human specimens \\
\bottomrule
\end{tabular}
\end{table}

\section{Methods}

\subsection{Cell lines and culture}

The CA1 oral squamous cell carcinoma cell line and its EMT-stem derivative sub-line were maintained under standard culture conditions as previously described \citep{Grosse-Wilde2015}. Normal oral keratinocytes and cancer-associated fibroblasts were isolated from surgical specimens and cultured as controls.

\subsection{Immunofluorescent staining}

Formalin-fixed paraffin-embedded (FFPE) sections were stained using a four-colour immunofluorescence protocol. The primary panel consisted of anti-EpCAM (yellow channel), anti-Vimentin (red channel), anti-CD24 (green channel), and DAPI (blue, nuclear counterstain). A control panel substituted anti-pan-keratin for anti-EpCAM. Antibody concentrations, incubation times, and antigen retrieval conditions were optimized for FFPE tissue.

\subsection{Automated imaging and analysis}

Entire histopathological slides were imaged using an automated fluorescence microscopy platform, generating over 12,000 imaging fields across 84 specimens. Image segmentation identified individual cells and quantified marker co-expression at the single-cell level. An EpCAM density map was computed across each slide to objectively delineate the tumour body (high EpCAM density regions) from the surrounding stroma (low EpCAM density).

\subsection{Tumour specimens}

Eighty-four human OSCC specimens were collected from surgical archives, stratified by metastatic status (presence or absence of lymph node metastasis at diagnosis). Specimens were divided into two independent batches for training and validation of the machine learning classifier. Ethical approval was obtained from the institutional review board.

\subsection{Machine learning classification}

An artificial neural network was trained on grey-level intensities of the three fluorescent markers (EpCAM, Vimentin, CD24) extracted from imaging tiles. Tiles were labeled by their tumour of origin (metastatic or non-metastatic). Ten-fold cross-validation was performed on batch 1, and the trained model was applied to batch 2 as an independent test set. Classification accuracy, sensitivity, and specificity were computed at each fold and averaged.

\subsection{Statistical analysis}

Cell line marker quantification was compared using two-tailed Student's $t$-tests with mean $\pm$ 95\% confidence intervals. Enrichment of triple-positive cells in metastatic versus non-metastatic stroma was assessed using appropriate significance tests. Machine learning performance was evaluated by 10-fold cross-validation accuracy and independent test set accuracy.

\section{Acknowledgments}

We thank the patients who consented to tissue banking and the histopathology teams for specimen preparation. This work was supported by [funding agencies].

\bibliographystyle{plainnat}
\bibliography{references}

\end{document}
